\documentclass[UTF8, scheme=chinese]{ctexart}

\usepackage{pdflscape}% 引入横向排版宏包

\usepackage[margin=0.5cm]{geometry}

\usepackage{xeCJK}
\xeCJKsetup{
  AutoFallBack = true,
  AutoFakeBold = true,
  PunctStyle = kaiming,
  CheckFullRight = true,
  CheckSingle = false,
  AllowBreakBetweenPuncts = true,
  CJKmath = true,
}

% 完整的中文章节格式设置
\ctexset{
  % 章设置
  % chapter = {
  % format = \huge\bfseries\centering,
  % nameformat = {},
  % titleformat = {},
  % name = {第,章},
  % number = \chinese{chapter},
  % aftername = \quad,
  % beforeskip = 20pt,
  % afterskip = 20pt,
  % },
  %   节设置
  section = {
    format = \Large\bfseries,
    nameformat = {},
    titleformat = {},
    name = {第,节},
    number = \chinese{section},
    aftername = \quad,
  },
  % 小节设置
  subsection = {
    format = \large\bfseries,
    nameformat = {},
    titleformat = {},
    name = {第,小节},
    number = \chinese{subsection},
    aftername = \quad,
  },
  % 子小节设置
  subsubsection = {
    format = \normalsize\bfseries,
    nameformat = {},
    titleformat = {},
    name = {第,子小节},
    number = \chinese{subsubsection},
    aftername = \quad,
  }
}

\xeCJKDeclareSubCJKBlock{Ext-A}{ "3400 -> "4DB5 }
\xeCJKDeclareSubCJKBlock{Ext-B}{ "20000 -> "2A6DF }
\xeCJKDeclareSubCJKBlock{CircNum}{ "2460 -> "2473 }
\xeCJKDeclareSubCJKBlock{Symbols}{ "2700 -> "27FF }

\setCJKmainfont[%
BoldFont={FZYouHei_GBK 513B},   %
% Ext-A={仓耳今楷05-W04},
Ext-B={SimSun-ExtB},
CircNum={IosevkaTerm NFM},
Symbols={IosevkaTerm NFM}
]
{NSimSun}
\newCJKfontfamily[li]\hmls{FZZJ-HMLSJF}
\newCJKfontfamily[kai]\fzxk{FZNewKai-Z03}
\newCJKfontfamily[song]\ssb{SimSun-ExtB}

\usepackage{xcolor}
\usepackage{pifont}

\usepackage{tikz}


% 引入颜色宏包（可选，用于让中文标注颜色变浅一点）
\usepackage{xpinyin}

% 记得将asmmath放在前面
\usepackage{amsmath}
% \usepackage{mathabx}
% \usepackage{fdsymbol}

% 更多功能的itemize以及表格
\usepackage{enumitem}

% 定义名为 descnum 的新环境
\newlist{descnum}{description}{1}
\setlist[descnum]{
  label={\arabic*.},
  font=\bfseries,
  leftmargin=2em,
  before=\stepcounter{section} % 可选：如果希望每章重新计数
}

\usepackage{tabularray}
\UseTblrLibrary{varwidth}
% 1. 修改当前页底部的提示文字 (默认是 Continued on next page)
\DefTblrTemplate{contfoot-text}{default}{下页继续}

% 2. 修改下一页顶部的提示文字 (默认是 (Continued))
\DefTblrTemplate{conthead-text}{default}{(续表)}

\usepackage{hyperref}
% 设置链接颜色，方便在 PDF 中看出哪里可以点击
\hypersetup{colorlinks=true, linkcolor=blue}

\title{\fzxk {\huge 两版《清诗300首》选诗对比}}
% \date{2026/04/11}

\begin{document}

\maketitle
% \tableofcontents

%% -*- coding: utf-8 -*-

\section{介绍}

目前可以看到有两本《清诗300首》，其一为\textbf{钱仲联、钱学增}于1983年初选，1992年定稿的版本，\textbf{叶嘉莹}为之序，所选清诗年代为公元1644年 ～ 1911年；其二为\textbf{刘世南}于1995年选定，范围亦包含有清一朝。另外，\textbf{钱仲联}还有《近代诗300首》，亦作为对比对比于表中。

出于个人兴趣，想看看两个选本的不同，因此做了下面的表格作为对照。

% 使用 landscape 环境包裹你的表格
% \begin{landscape}
%   \pagestyle{empty}
\begin{center}
  \begin{longtblr}[
    label = none,
    ]{
    vlines, hlines,
    % 移除全局的 halign=c，改为在每一列单独设置，这样更灵活
    columns={valign=m},
    column{1}={0.1\linewidth, halign=c, bg=azure9, font=\large\fzxk},
    % column{2}={0.1\linewidth, halign=c},
    % 关键修改：给第三列指定宽度（X代表填满剩余所有空间），并设为左对齐
    column{2,3,4}={0.267\linewidth, halign=l},
    rows={abovesep=5pt, belowsep=5pt},
    row{1}={font=\Large\bfseries, bg=azure7, halign=c}, % 表头居中
    % 关键修改：处理单元格内的 list 环境
    measure=vbox,
    stretch=-1,
    }
    作者                       & 刘选本                                             & 钱选本                                             & 钱选近代诗                                             \\
    钱谦益                     &                                                    & 古诗赠新城王贻上                                   &                                                        \\
                               &                                                    & 吴门春仲送李生还长干                               &                                                        \\
                               &                                                    & 西湖杂感（二十首选六）                             &                                                        \\
                               &                                                    & 金陵秋兴八首次草堂韵己亥七月初一作（选二）         &                                                        \\
                               &                                                    & 迎神曲（十二首选四）                               &                                                        \\
                               & 陆宣公墓道行                                       &                                                    &                                                        \\
                               & 简侯研德，并示记原                                 &                                                    &                                                        \\
                               & 题沈朗倩石崖秋柳小景                               &                                                    &                                                        \\
    金人瑞                     & 与儿子                                             &                                                    &                                                        \\
    冯班                       &                                                    & 有赠                                               &                                                        \\
                               &                                                    & 紫媛（自北归吴）                                   &                                                        \\
    吴伟业                     &                                                    & 遇南厢园叟感赋八十韵                               &                                                        \\
                               &                                                    & 清凉山赞佛诗（四首）                               &                                                        \\
                               &                                                    & 圆圆曲                                             &                                                        \\
                               &                                                    & 打冰词                                             &                                                        \\
                               &                                                    & 过吴江有感                                         &                                                        \\
                               &                                                    & 过淮阴有感（二首选一）                             &                                                        \\
                               &                                                    & 下相极乐庵读同年北使时诗卷                         &                                                        \\
                               &                                                    & 阻雪                                               &                                                        \\
                               & 永和宫词                                           &                                                    &                                                        \\
                               & 梅村                                               &                                                    &                                                        \\
                               & 采石矶                                             &                                                    &                                                        \\
    方以智                     &                                                    & 独往                                               &                                                        \\
    杜濬                       &                                                    & 楼夕                                               &                                                        \\
                               & 泰州绝句                                           &                                                    &                                                        \\
                               & 闵友谈水东之胜                                     &                                                    &                                                        \\
    钱秉镫                     &                                                    & 扬州访汪辰初                                       &                                                        \\
    \hyperlink{XingFang}{邢昉} & 印度在禅河庵，与予寺舍相距五里，纪事（选一）       &                                                    &                                                        \\
                               & 赠祖心上人                                         &                                                    &                                                        \\
    钱澄之                     & 田园杂诗（选一）                                   &                                                    &                                                        \\
                               & 别孙不害，途中有怀                                 &                                                    &                                                        \\
    顾炎武                     & 赴东（六首，并序）                                 &                                                    &                                                        \\
                               & 楼观                                               &                                                    &                                                        \\
                               &                                                    & 秋山二首                                           &                                                        \\
                               &                                                    & 白下                                               &                                                        \\
                               &                                                    & 汾州祭吴炎潘柽章二节士                             &                                                        \\
                               &                                                    & 雨中至华下宿王山史家                               &                                                        \\
    宋琬                       &                                                    & 同欧阳令饮凤凰山下（选一）                         &                                                        \\
                               & ✓                                                  & 舟中见猎犬有感                                     &                                                        \\
                               & 剡杂溪道中                                         &                                                    &                                                        \\
                               & 检阅故人姜筼筜遗稿，泫然有作                       &                                                    &                                                        \\
    龚鼎孳                     & ✓                                                  & 上巳将过金陵（四首选一）                           &                                                        \\
    尤侗                       &                                                    & 煮粥行                                             &                                                        \\
    申涵光                     & 春雪歌                                             &                                                    &                                                        \\
                               & 答人                                               &                                                    &                                                        \\
    吴嘉纪                     &                                                    & 哀羊裘为孙八赋                                     &                                                        \\
                               &                                                    & 一钱行赠林茂之                                     &                                                        \\
                               &                                                    & 李家娘                                             &                                                        \\
                               &                                                    & 打鲥鱼（二首）                                     &                                                        \\
                               &                                                    & 九日寄徐式家                                       &                                                        \\
                               &                                                    & 内人生日                                           &                                                        \\
                               &                                                    & 玉钩斜                                             &                                                        \\
                               & 我昔（三首选一）                                   &                                                    &                                                        \\
                               & 夜发                                               &                                                    &                                                        \\
                               & 绝句（白头灶户）                                   &                                                    &                                                        \\
    施闰章                     & 临江悯旱                                           &                                                    &                                                        \\
                               & 山行                                               &                                                    &                                                        \\
                               &                                                    & 浮萍兔丝篇                                         &                                                        \\
                               &                                                    & 过湖北山家                                         &                                                        \\
                               &                                                    & 燕子矶                                             &                                                        \\
                               &                                                    & 舟中立秋                                           &                                                        \\
                               & ✓                                                  & 泊樵舍                                             &                                                        \\
                               &                                                    & 至南旺                                             &                                                        \\
    洪昇                       & 将游大梁                                           &                                                    &                                                        \\
    毛奇龄                     &                                                    & 打虎儿行                                           &                                                        \\
    汪琬                       &                                                    & 官军行                                             &                                                        \\
                               & 东归道中（三首选一）                               &                                                    &                                                        \\
    陈维崧                     & 刘逸民隐如                                         &                                                    &                                                        \\
                               & 春暮杂忆，和彭羡门韵                               &                                                    &                                                        \\
                               &                                                    & 将归留别练塘诸子                                   &                                                        \\
    蒋超                       &                                                    & 金陵旧院                                           &                                                        \\
    缪彤                       &                                                    & 渡江                                               &                                                        \\
    钱曾                       &                                                    & 秋夜宿破山寺绝句（十二首选一）                     &                                                        \\
    朱尊彝                     & 大孤山                                             &                                                    &                                                        \\
                               & 晚次崞县                                           &                                                    &                                                        \\
                               & 为钱给事（晋锡）题王给事（原祁）富春大岭图（二首） &                                                    &                                                        \\
                               &                                                    & 马草行                                             &                                                        \\
    屈大均                     & 赠朱士稚                                           &                                                    &                                                        \\
                               & 王不庵作卧龙松歌为余寿，诗以酬之                   &                                                    &                                                        \\
                               & 摄山秋夕作                                         &                                                    &                                                        \\
                               & 人日衡阳道中                                       &                                                    &                                                        \\
                               & 紫蒙                                               &                                                    &                                                        \\
                               &                                                    & 登罗浮绝顶奉同蒋王二大夫作                         &                                                        \\
                               &                                                    & 读陈胜传                                           &                                                        \\
                               &                                                    & 秣陵                                               &                                                        \\
                               &                                                    & 云州秋望                                           &                                                        \\
                               &                                                    & 江皋                                               &                                                        \\
                               &                                                    & 天边                                               &                                                        \\
    陈恭尹                     &                                                    & 发舟寄湛用喈钟裴仙湛天石                           &                                                        \\
                               &                                                    & 崖门谒三忠祠                                       &                                                        \\
                               & 秋夜同梁巨川、卢偶斯过宿潘子登寓斋……次子登韵       &                                                    &                                                        \\
    吴兆骞                     &                                                    & 奉送巴大将军东征逻察                               &                                                        \\
                               &                                                    & 秋夜师次松花江大将军……即成口号示同观诸子           &                                                        \\
                               &                                                    & 黑林                                               &                                                        \\
                               & 榆关老翁行                                         &                                                    &                                                        \\
    陆次云                     &                                                    & 出门（二首）                                       &                                                        \\
                               &                                                    & 咏史                                               &                                                        \\
    王士禛                     &                                                    & 定军山诸葛公墓下作                                 &                                                        \\
                               &                                                    & 息斋夜宿即景有怀故园                               &                                                        \\
                               &                                                    & 雨后观音门渡江                                     &                                                        \\
                               &                                                    & 晓雨重登燕子矶绝顶作                               &                                                        \\
                               &                                                    & 再过露筋祠                                         &                                                        \\
                               &                                                    & 秦淮杂诗（十四首选五）                             &                                                        \\
                               &                                                    & 冶春绝句（二十首选三）                             &                                                        \\
                               &                                                    & 蟂矶灵泽夫人祠（二首选一）                         &                                                        \\
                               & 余澹心寄金陵咏怀古迹诗，却寄（二手选一）           &                                                    &                                                        \\
                               & 青山                                               &                                                    &                                                        \\
                               & 送陈其年归宜兴（二首）                             &                                                    &                                                        \\
                               & 雨中度故关                                         &                                                    &                                                        \\
    汤右曾                     & 泊长沙西门，追念旧游有感                           &                                                    &                                                        \\
    宋荦                       &                                                    & 海上杂诗（三首选一）                               &                                                        \\
    许虬                       &                                                    & 折杨柳歌                                           &                                                        \\
    吴雯                       &                                                    & 明妃                                               &                                                        \\
    查慎行                     &                                                    & 麻阳运船行                                         &                                                        \\
                               &                                                    & 初入黔境土人皆居悬岩峭壁间……心恻而作是诗           &                                                        \\
                               &                                                    & 秦邮道中即目                                       &                                                        \\
                               &                                                    & 邺下杂咏（四首选一）                               &                                                        \\
                               & 拂水山庄                                           &                                                    &                                                        \\
                               & 残冬展假，病榻消寒，……录存十六章（选一）           &                                                    &                                                        \\
    纳兰性德                   &                                                    & 秣陵怀古                                           &                                                        \\
    梁佩兰                     & 独立                                               &                                                    &                                                        \\
    赵执信                     & 行十八滩中                                         &                                                    &                                                        \\
                               & 出都                                               &                                                    &                                                        \\
                               & 昭阳湖行，书所见                                   &                                                    &                                                        \\
                               &                                                    & 甿入城行                                           &                                                        \\
    沈德潜                     &                                                    & 制府来                                             &                                                        \\
                               & 凿冰行                                             &                                                    &                                                        \\
                               & 后凿冰行                                           &                                                    &                                                        \\
                               & 过许州                                             &                                                    &                                                        \\
    黄任                       &                                                    & 赈粥行                                             &                                                        \\
                               & ✓                                                  & 西湖杂诗（十四首选三）                             &                                                        \\
                               & 归舟杂诗（二十六首选一）                           &                                                    &                                                        \\
    厉鹗                       &                                                    & 晚登韬光绝顶                                       &                                                        \\
                               &                                                    & 开平王孙种菜歌并序                                 &                                                        \\
                               &                                                    & 灵隐寺月夜                                         &                                                        \\
                               &                                                    & 入秋酷暑七月十二日晚风月甚清遂有凉意               &                                                        \\
                               & ✓                                                  & 湖楼题壁                                           &                                                        \\
                               &                                                    & 自石湖至横塘                                       &                                                        \\
                               &                                                    & 归舟江行望燕子矶作                                 &                                                        \\
                               & 次韵堇浦春日书怀（四首选一）                       &                                                    &                                                        \\
                               & 春寒                                               &                                                    &                                                        \\
                               & 西湖春雨，四首（选一）                             &                                                    &                                                        \\
    杭世骏                     & 出国门作（二首选一）                               &                                                    &                                                        \\
                               & 沂州道中，忆故园梅信                               &                                                    &                                                        \\
    姚范                       & 书家书后                                           &                                                    &                                                        \\
    徐兰                       &                                                    & 出关                                               &                                                        \\
    郑燮                       &                                                    & 还家行                                             &                                                        \\
                               &                                                    & 私刑恶                                             &                                                        \\
                               & 悍吏                                               &                                                    &                                                        \\
    严遂成                     &                                                    & 砥柱峰                                             &                                                        \\
                               &                                                    & 白水岩瀑布                                         &                                                        \\
                               &                                                    & 渡鄱阳                                             &                                                        \\
                               &                                                    & 龙泉关                                             &                                                        \\
                               &                                                    & 宿许天植见山楼                                     &                                                        \\
                               &                                                    & 三垂冈                                             &                                                        \\
                               &                                                    & 秋夜投止山家                                       &                                                        \\
                               &                                                    & 安肃道中                                           &                                                        \\
    桑调元                     &                                                    & 五人墓                                             &                                                        \\
    胡天游                     &                                                    & 烈女李三行                                         &                                                        \\
                               &                                                    & 晓行                                               &                                                        \\
                               & 万泉叹                                             &                                                    &                                                        \\
                               & 塞上观落日                                         &                                                    &                                                        \\
    倪瑞璿                     &                                                    & 闻蛙                                               &                                                        \\
    王又曾                     &                                                    & 登报恩寺浮图                                       &                                                        \\
                               &                                                    & 桐江                                               &                                                        \\
                               &                                                    & 经天姥寺                                           &                                                        \\
                               &                                                    & 临平道中看荷花同朱冰壑陈渔沂（二首）               &                                                        \\
    钱载                       &                                                    & 清远道士养鹤涧                                     &                                                        \\
                               &                                                    & 到家作（四首选二）                                 &                                                        \\
                               &                                                    & 清流关                                             &                                                        \\
                               &                                                    & 望石湖                                             &                                                        \\
                               &                                                    & 出东林六七里望庐山                                 &                                                        \\
                               & 检先孺人遗箧，得（载）己亥康熙五十八年雪夜诗       &                                                    &                                                        \\
    金农                       & 杨三十五过广陵客舍                                 &                                                    &                                                        \\
                               & 瘦马                                               &                                                    &                                                        \\
    袁枚                       &                                                    & 行十里至黄崖再登文殊塔观瀑                         &                                                        \\
                               &                                                    & 澶渊                                               &                                                        \\
                               &                                                    & 秦中杂感（八首选一）                               &                                                        \\
                               & ✓                                                  & 鸡                                                 &                                                        \\
                               &                                                    & 马嵬（四首选一）                                   &                                                        \\
                               & 十一月十三日，冷水步夜起玩月                       &                                                    &                                                        \\
                               & 所见                                               &                                                    &                                                        \\
                               & 春日杂诗（二首）                                   &                                                    &                                                        \\
    王鸣盛                     & 西湖葛岭有嘲                                       &                                                    &                                                        \\
    钱大昕                     & 读曲歌（二首选一）                                 &                                                    &                                                        \\
                               & 即事                                               &                                                    &                                                        \\
    蒋士铨                     & 哭杨子载（二首选一）                               &                                                    &                                                        \\
    李宪噩                     & 雪后归南村，道中作                                 &                                                    &                                                        \\
    曾燠                       & 山烧                                               &                                                    &                                                        \\
    阮元                       & 莱州蜉蝣岛                                         &                                                    &                                                        \\
    冯廷櫆                     & 郎陵行                                             &                                                    &                                                        \\
    赵翼                       & 古诗二十首（选一）                                 &                                                    &                                                        \\
                               & 归田即事（四首选一）                               &                                                    &                                                        \\
                               &                                                    & 杂题八首（选一）                                   &                                                        \\
                               &                                                    & 闲居读书（六首选一）                               &                                                        \\
    姚鼐                       &                                                    & 岳州城上                                           &                                                        \\
                               & 钱舜举《萧翼赚兰亭图》                             &                                                    &                                                        \\
                               & 登永济寺阁，寺是中山王旧园                         &                                                    &                                                        \\
                               & 金陵晓发                                           &                                                    &                                                        \\
    翁方纲                     &                                                    & 望罗浮                                             &                                                        \\
    孙星衍                     & 夜起步月，偕妇王采薇                               &                                                    &                                                        \\
    洪亮吉                     & 自勉（二首）                                       &                                                    &                                                        \\
                               & 晓行                                               &                                                    &                                                        \\
                               & 悯旱                                               &                                                    &                                                        \\
                               &                                                    & 松树塘万松歌                                       &                                                        \\
    吴锡麒                     &                                                    & 狮子林歌                                           &                                                        \\
                               &                                                    & 韬光庵                                             &                                                        \\
                               & 读放翁集                                           &                                                    &                                                        \\
                               & 虎邱                                               &                                                    &                                                        \\
    黎简                       &                                                    & 龙门滩                                             &                                                        \\
                               &                                                    & 亡妹生日                                           &                                                        \\
                               &                                                    & 村饮                                               &                                                        \\
                               & 邕州城楼                                           &                                                    &                                                        \\
    黄景仁                     &                                                    & 圈虎行                                             &                                                        \\
                               & ✓                                                  & 金陵别邵大仲游                                     &                                                        \\
                               &                                                    & 都门秋思（四首选一）                               &                                                        \\
                               &                                                    & 新安滩                                             &                                                        \\
                               &                                                    & 幼女                                               &                                                        \\
                               &                                                    & 山馆夜作                                           &                                                        \\
                               &                                                    & 别老母                                             &                                                        \\
                               & 得稚存、渊如书，却寄                               &                                                    &                                                        \\
                               & 新凉曲                                             &                                                    &                                                        \\
    宋湘                       &                                                    & 灌花吟                                             &                                                        \\
                               &                                                    & 买鱼歌                                             &                                                        \\
                               &                                                    & 河南道中书事感怀（五首选二）                       &                                                        \\
                               &                                                    & 入洞庭                                             &                                                        \\
                               &                                                    & 贵州飞云洞题壁                                     &                                                        \\
                               &                                                    & 祭风台                                             &                                                        \\
    孙原湘                     & 澄海楼望海                                         &                                                    &                                                        \\
                               & 月午楼歌                                           &                                                    &                                                        \\
                               &                                                    & 征妇怨                                             &                                                        \\
                               &                                                    & 蒙山                                               &                                                        \\
                               &                                                    & 西陵峡                                             &                                                        \\
    吴震                       &                                                    & 月夜游孤山                                         &                                                        \\
    王昙                       &                                                    & 善才生二十五月矣计识得二百五十余字示以诗云         &                                                        \\
                               &                                                    & 焦山夜泊                                           &                                                        \\
    张问陶                     & ✓                                                  & 十六夜雪中渡江                                     &                                                        \\
                               &                                                    & 出栈（二首选一）                                   &                                                        \\
                               &                                                    & 阳湖道中                                           &                                                        \\
                               & 摇落                                               &                                                    &                                                        \\
    屠倬                       & 题友人斋壁                                         &                                                    &                                                        \\
                               & 送频伽还魏塘                                       &                                                    &                                                        \\
    姚莹                       & 三十                                               &                                                    &                                                        \\
                               & 吴桥暮雨                                           &                                                    &                                                        \\
    舒位                       &                                                    & 卧闻蟋蟀偶成                                       &                                                        \\
                               &                                                    & 杨花诗                                             &                                                        \\
                               &                                                    & 读《文选》诗九首（选一）                           &                                                        \\
                               &                                                    & 花生日诗魏塘道中作                                 &                                                        \\
                               & 送客至莫愁湖上，遂登胜棋楼有作（二首选一）         &                                                    &                                                        \\
    吴嵩梁                     & 王荆公祠                                           &                                                    &                                                        \\
    郭麐                       & 秋感二首（选一）                                   &                                                    &                                                        \\
                               & 阿桐生日                                           &                                                    &                                                        \\
                               & 《杨、陆谈诗图》，为退庵题                         &                                                    &                                                        \\
                               & 喜家书至                                           &                                                    &                                                        \\
                               & 绝句（沉香火冷）                                   &                                                    &                                                        \\
                               &                                                    & 新晴即事（六首选一）                               &                                                        \\
    彭兆荪                     & 独坐                                               &                                                    &                                                        \\
                               & 滁阳                                               &                                                    &                                                        \\
    韩是升                     & 梦谒顾亭林先生墓，得句云……遂足成之                 &                                                    &                                                        \\
    陈寿祺                     & 春暮意行近郊，偶赋                                 &                                                    &                                                        \\
    陈文述                     &                                                    & 夏日杂诗（七首选一）                               &                                                        \\
                               & 金陵杂感（二首选一）                               &                                                    &                                                        \\
    邓显鹤                     & 过襄阳，书寄家遁溪丈                               &                                                    &                                                        \\
    张维屏                     &                                                    & 侠客行                                             &                                                        \\
                               &                                                    & 三元里                                             & ✓                                                      \\
                               &                                                    & 秋霖                                               &                                                        \\
    钱仪吉                     &                                                    & 读黎二樵诗                                         &                                                        \\
                               &                                                    & 山行                                               &                                                        \\
                               & 樾亭弟至京夜话（四首选二）                         &                                                    &                                                        \\
    陈沆                       &                                                    & 项师竹张馥亭自麻城来访欣然有作                     &                                                        \\
                               &                                                    & 河南道上乐府四章（选二）                           &                                                        \\
                               &                                                    & 夜抵刘山人家                                       &                                                        \\
                               &                                                    & 扬州城楼                                           &                                                        \\
                               &                                                    & 渡河遇相识寄家书                                   &                                                        \\
    潘德舆                     &                                                    & 鹊巢                                               &                                                        \\
    龚自珍                     & 寒月吟（五首选一）                                 &                                                    &                                                        \\
                               & 暮雨谣三叠                                         &                                                    &                                                        \\
                               &                                                    & 自春徂秋偶有所触拉杂书之漫不诠次得十五首（选二）   &                                                        \\
                               &                                                    & 能令公少年行                                       &                                                        \\
                               &                                                    & 咏史                                               &                                                        \\
                               &                                                    & 梦中作四截句（选一）                               &                                                        \\
                               &                                                    & 歌筵有乞书扇者                                     &                                                        \\
                               &                                                    & 猛忆                                               &                                                        \\
                               & ✓                                                  & 己亥杂诗（三百十五首选二）                         & ✓                                                      \\
                               &                                                    &                                                    & 伪鼎行                                                 \\
                               &                                                    &                                                    & 夜坐(二首选一)                                         \\
                               &                                                    &                                                    & 释言四首之一                                           \\
    魏源                       &                                                    & 晓窗                                               &                                                        \\
                               &                                                    & 天台赠僧（二首选一）                               &                                                        \\
                               & 重游盘山寺                                         &                                                    &                                                        \\
                               & 高邮州署，秋日偶题（三首）                         &                                                    &                                                        \\
                               &                                                    &                                                    & 天台石梁雨后观瀑歌                                     \\
    何绍基                     &                                                    & 飞云岩                                             &                                                        \\
                               &                                                    & 元象                                               &                                                        \\
                               & 同子毅弟早起至岱顶                                 &                                                    &                                                        \\
                               & 山人诗舲柬至，谓名酒如美人……口占奉谢               &                                                    &                                                        \\
                               & 逆风                                               &                                                    &                                                        \\
                               & 补竹（选二）                                       &                                                    &                                                        \\
                               &                                                    &                                                    & 题陈忠愍公遗像练栗人属作                               \\
                               &                                                    &                                                    & 晓发看月                                               \\
    汤鹏                       & 南邻                                               &                                                    &                                                        \\
    汪士铎                     & 蟋蟀                                               &                                                    &                                                        \\
                               & 挽林文忠                                           &                                                    &                                                        \\
    朱琦                       &                                                    & 王刚节公家传书后                                   &                                                        \\
                               & 当阳道中                                           &                                                    &                                                        \\
                               &                                                    &                                                    & 老兵叹                                                 \\
    鲁一同                     &                                                    & (辛丑)重有感（八首选一）                           & ✓                                                      \\
                               & 明月                                               &                                                    &                                                        \\
                               & 忆长安旧游（二首选一）                             &                                                    &                                                        \\
                               &                                                    &                                                    & 读史杂感(五首选一)                                     \\
    朱次琦                     & 读史（二首）                                       &                                                    &                                                        \\
    佘梅                       &                                                    & 听屠生说马僧事证之随园所书者纪以古诗               &                                                        \\
    姚燮                       & 夷人退还定海后，纪诗四章（选一）                   &                                                    &                                                        \\
                               & 闻定海陷五章（选二）                               &                                                    &                                                        \\
                               & 当古谣十三章（选二）                               &                                                    &                                                        \\
                               &                                                    & 北村妇                                             &                                                        \\
                               &                                                    & 双鸩篇                                             &                                                        \\
                               &                                                    &                                                    & 由妙庄严路至普济寺三章（选一）                         \\
                               &                                                    &                                                    & 笑天狮子岭                                             \\
                               &                                                    &                                                    & 法华洞                                                 \\
                               &                                                    &                                                    & 孤凄坑                                                 \\
                               &                                                    &                                                    & 沈家庄夜半大雨                                         \\
                               &                                                    &                                                    & 兵巡街                                                 \\
                               &                                                    &                                                    & 自华桃横冒云下紫树陵                                   \\
                               &                                                    &                                                    & 澄灵涧                                                 \\
    郑献甫                     & 杂述（选一）                                       &                                                    &                                                        \\
    郑珍                       & 黄焦石                                             & ✓                                                  &                                                        \\
                               & 赠刘生子莹（之琇），生布衣，业烛笼为活             &                                                    &                                                        \\
                               & 正月十六戏书                                       &                                                    &                                                        \\
                               & 贵阳寄内                                           &                                                    &                                                        \\
                               & 三坡曲（二首）                                     &                                                    &                                                        \\
                               &                                                    & 候涨退                                             &                                                        \\
                               &                                                    & 春尽日                                             &                                                        \\
                               &                                                    & 云门墱                                             &                                                        \\
                               &                                                    & 怀阳洞                                             &                                                        \\
                               &                                                    & 江边老叟诗                                         &                                                        \\
                               & ✓                                                  & 自沾益出宣威入东川                                 &                                                        \\
                               &                                                    & 题仇实父《清明上河图》                             &                                                        \\
                               &                                                    &                                                    & 望乡吟                                                 \\
                               &                                                    &                                                    & 白水瀑布                                               \\
                               &                                                    &                                                    & 自毛口宿花堌                                           \\
                               &                                                    &                                                    & 次杨林晚望                                             \\
                               &                                                    &                                                    & 邯郸                                                   \\
    邵懿辰                     & 偶忆（五首选一）                                   &                                                    &                                                        \\
                               & 喜闻少鹤病已                                       &                                                    &                                                        \\
    贝青乔                     &                                                    & 军中杂诔诗（十八首选一）                           &                                                        \\
                               & ✓                                                  & 咄咄吟（一百二十四首选一）                         & ✓                                                      \\
                               &                                                    & 自编军中纪事诗《咄咄吟》朋旧多题赠之为答（五选一） &                                                        \\
                               & 赤津岭                                             &                                                    & ✓                                                      \\
                               & 自临安至於潜，夜宿浮溪旅店作                       &                                                    &                                                        \\
    曾国藩                     &                                                    & 傲奴                                               &                                                        \\
                               & 送梅伯言归金陵                                     &                                                    &                                                        \\
                               & 读李义山诗集                                       &                                                    &                                                        \\
                               &                                                    &                                                    & 寄弟(三首选一)                                         \\
    戴钧衡                     & 得张亨甫开封书                                     &                                                    &                                                        \\
    周寿昌                     & 谒于忠肃公墓                                       &                                                    &                                                        \\
                               & 南康野次书所见                                     &                                                    &                                                        \\
                               & 秋夜独坐                                           &                                                    &                                                        \\
    刘蓉                       & 园花为风雨摧落，感赋长句                           &                                                    &                                                        \\
    莫友芝                     &                                                    & 霸王坡                                             &                                                        \\
                               & 公安县                                             &                                                    &                                                        \\
                               & 草堂杂诗，三首（选一）                             &                                                    &                                                        \\
                               & 神洲滩谣二首（选一）                               &                                                    &                                                        \\
                               & 横石                                               &                                                    &                                                        \\
    沈谨学                     & 春日杂诗（二首选一）                               &                                                    &                                                        \\
                               & 溪上                                               &                                                    &                                                        \\
                               & 吴淞江归棹                                         &                                                    &                                                        \\
    张裕钊                     & 无题（二首）                                       &                                                    &                                                        \\
    张祖继                     & 春雨                                               &                                                    &                                                        \\
    杨文莹                     & 涧石（二首）                                       &                                                    &                                                        \\
                               & 示儿，二绝（选一）                                 &                                                    &                                                        \\
    陈豪                       & 读史记，感齐姜及赵衰妻事得两绝句                   &                                                    &                                                        \\
                               &                                                    &                                                    & 题画(二首选一)                                         \\
    王存原                     & 丁巳八月十三日过菜市，吊晚翠（三首选二）           &                                                    &                                                        \\
    秦树生                     & 为人题画兰（三首）                                 &                                                    &                                                        \\
    孔宪彝                     & 阳雨谣                                             &                                                    &                                                        \\
    薛时雨                     & 湖州道中（二首选一）                               &                                                    &                                                        \\
    郭嵩焘                     & 和朱暝庵《岁暮杂感》（三首选一）                   &                                                    &                                                        \\
                               & 九日                                               &                                                    &                                                        \\
                               & 平原道中                                           &                                                    &                                                        \\
    江湜                       & 秋感（二首选一）                                   &                                                    &                                                        \\
                               & 将抵邵武                                           &                                                    &                                                        \\
                               & 途中书意                                           &                                                    &                                                        \\
                               & 近年                                               &                                                    &                                                        \\
                               & 岁除日戏作二诗                                     &                                                    &                                                        \\
                               & 舟中二绝（选一）                                   &                                                    & ✓                                                      \\
                               &                                                    & 湖楼早起（二首）                                   &                                                        \\
                               &                                                    &                                                    & 道中忆旧仆沈用作四诗以酬昔劳(选二)                     \\
                               &                                                    &                                                    & 一风                                                   \\
                               &                                                    &                                                    & 彦冲画柳燕                                             \\
                               &                                                    &                                                    & 连夜月下有作                                           \\
                               &                                                    &                                                    & 南台酒家题壁                                           \\
                               &                                                    &                                                    & 野兴(六首选一)                                         \\
                               &                                                    &                                                    & 雨中                                                   \\
    洪仁玕                     &                                                    &                                                    & 回港舟中诗                                             \\
                               &                                                    &                                                    & 二月下浣，军次遂安城北，吟于行府                       \\
    姚濬昌                     & 宿范水即事，是去年经乱处                           &                                                    &                                                        \\
    易佩绅                     & 湖因水涸而曲，略写其状                             &                                                    &                                                        \\
    金和                       &                                                    & 兰陵女儿行                                         &                                                        \\
                               & 初七日去大营，拟寄城中诸友                         &                                                    &                                                        \\
                               &                                                    &                                                    & 断指生歌                                               \\
    蒋春霖                     & 冬柳                                               &                                                    &                                                        \\
    符兆伦                     & 忆家                                               &                                                    &                                                        \\
    蒋湘南                     & 醉来                                               &                                                    &                                                        \\
    戴家麟                     & 三汊河晚宿                                         &                                                    &                                                        \\
    周星誉                     & 春日西湖即席                                       &                                                    &                                                        \\
    邓辅纶                     & 登衡山南天门                                       &                                                    &                                                        \\
                               &                                                    & 鸿雁篇（三首）                                     & ✓                                                      \\
                               &                                                    & 听雨轩坐秋                                         &                                                        \\
    翁同龢                     &                                                    & 甲辰五月二十日绝笔                                 &                                                        \\
                               & 病榻不寐                                           &                                                    &                                                        \\
                               & 衡门                                               &                                                    &                                                        \\
                               &                                                    &                                                    & 游西山见宝竹坡题名，因书其后                           \\
                               & 将之江右视筱珊侄                                   &                                                    &                                                        \\
    高心夔                     & 城西                                               &                                                    &                                                        \\
                               & 寄同年龚济南（易图）                               &                                                    &                                                        \\
                               &                                                    &                                                    & 鄱阳翁                                                 \\
    李慈铭                     &                                                    & 庚午书事二首                                       &                                                        \\
                               & 雨泊凤凰山下石泉庵                                 &                                                    &                                                        \\
                               & 和珊士送春诗，心字韵                               &                                                    &                                                        \\
                               & 甲申三月十三日出都，……绝句五十首（选二）           &                                                    &                                                        \\
                               &                                                    &                                                    & 边闻三首(选二)                                         \\
                               &                                                    &                                                    & 为潘星斋侍郎题秦宜亭岁晚江村小景三绝句(选二)           \\
    周星诒                     & 夜行萧山道中，望姚大舟不至                         &                                                    &                                                        \\
    王闿运                     & 泰山诗，孟冬朔日登山作                             &                                                    &                                                        \\
                               & 齐河道中雪行，偶作（二首选一）                     &                                                    &                                                        \\
                               &                                                    & 圆明园词                                           & ✓                                                      \\
                               &                                                    & 晓上空泠峡                                         &                                                        \\
                               &                                                    &                                                    & 独游妙相庵观道、咸诸卿相刻石                           \\
    张佩纶                     &                                                    &                                                    & 雁(五首选四)                                           \\
                               &                                                    &                                                    & 居庸                                                   \\
    张之洞                     &                                                    & 鸡鸣寺                                             &                                                        \\
                               &                                                    & 四月下旬过崇效寺访牡丹花已残损                     &                                                        \\
                               & 海水（二首选一）                                   &                                                    &                                                        \\
                               & 赠日本长冈护美                                     &                                                    &                                                        \\
                               & 金陵杂诗（十六首选一）                             &                                                    &                                                        \\
                               &                                                    &                                                    & 故府                                                   \\
    冯煦                       &                                                    &                                                    & 八月二十一日之夜，仆卧已久……亦成三首（选一）           \\
                               &                                                    &                                                    & 强漱泉也(选一)                                         \\
                               &                                                    &                                                    & 瓜步                                                   \\
                               &                                                    &                                                    & 送研孙归湘中(二首)                                     \\
    张洵                       & 书陆放翁《南园记》后（二首）                       &                                                    &                                                        \\
    丁尧臣                     & 阿房                                               &                                                    &                                                        \\
    曾纪泽                     & 八月十五日夜，森比德堡对月                         &                                                    &                                                        \\
                               & 香严书询近况，诗以代柬                             &                                                    &                                                        \\
    吴汝纶                     & 北征别张廉卿，即送其东游（二选一）                 &                                                    &                                                        \\
    袁昶                       & 北游诗五章（选一）                                 &                                                    &                                                        \\
                               & 和友人吴山水仙王祠                                 &                                                    &                                                        \\
                               & 寄榆园逸叟（二首选一）                             &                                                    &                                                        \\
                               &                                                    & 直房小憩                                           &                                                        \\
                               &                                                    &                                                    & 晓发                                                   \\
                               &                                                    &                                                    & 西轩睡起偶成绝句                                       \\
    叶大庄                     & 齐河大风，晚始得渡                                 &                                                    &                                                        \\
                               &                                                    &                                                    & 吴江舟中(三首选一)                                     \\
                               &                                                    &                                                    & 竹亭                                                   \\
                               &                                                    &                                                    & 村居书事四首(选二)                                     \\
                               &                                                    &                                                    & 往游洪塘沿路即目三首(选二)                             \\
                               &                                                    &                                                    & 溪堂闲居六首(选一)                                     \\
                               &                                                    &                                                    & 桐江舟中漫成寄兰溪补堂明府                             \\
    郭曾炘                     & 偕朗溪、芦台访熙民叔侄村居                         &                                                    &                                                        \\
    樊增祥                     & ✓                                                  & 八月六日过灞桥口占                                 &                                                        \\
                               & 雨中度棋盘岭                                       &                                                    &                                                        \\
                               & 闻都门消息（五首选一）                             &                                                    & ✓                                                      \\
                               & 东溪诗（并序二首）                                 &                                                    &                                                        \\
                               & 听竹                                               &                                                    &                                                        \\
                               & 奉和《金陵杂诗》十六首（选一）                     &                                                    &                                                        \\
    黄遵宪                     & 登巴黎铁塔                                         & ✓                                                  &                                                        \\
                               & 寄怀左子兴领事秉隆                                 &                                                    &                                                        \\
                               & 海行杂感（五首选二）                               &                                                    &                                                        \\
                               &                                                    & 今别离                                             &                                                        \\
                               &                                                    & 度辽将蒋春霖军歌                                   &                                                        \\
                               &                                                    & 雁                                                 &                                                        \\
                               &                                                    & 长沙吊贾谊宅                                       &                                                        \\
                               &                                                    &                                                    & 锡兰岛卧佛(六首选一)                                   \\
                               &                                                    &                                                    & 冯将军歌                                               \\
                               &                                                    &                                                    & 香港感怀十首(选一)                                     \\
                               &                                                    &                                                    & 酬曾重伯编修(二首选一)                                 \\
    陈宝琛                     &                                                    & 大悲寺秋海棠                                       &                                                        \\
                               & 山居                                               &                                                    &                                                        \\
                               &                                                    &                                                    & 缅侨叹                                                 \\
                               &                                                    &                                                    & 华严精舍(五首选二)                                     \\
                               &                                                    &                                                    & 感春                                                   \\
                               &                                                    &                                                    & 沧趣楼杂诗(九首选一)                                   \\
                               &                                                    &                                                    & 吉隆车中口号                                           \\
                               &                                                    &                                                    & 馆故甲必丹叶来宅……时有演说革命者，援此晓之             \\
    杨深秀                     &                                                    &                                                    & 仿元遗山论诗绝句五十首(选四)                           \\
    王甲荣                     &                                                    &                                                    & 彩云曲                                                 \\
    盛昱                       &                                                    &                                                    & 题钱南园画马一枯树一瘦马一小马                         \\
    沈曾植                     &                                                    & 题赵吴兴《鸥波亭图》                               &                                                        \\
                               &                                                    & 偕石遗渡江                                         &                                                        \\
                               &                                                    & 阁夜示证刚                                         &                                                        \\
                               &                                                    & 西湖杂诗（十七首选四）                             &                                                        \\
                               &                                                    &                                                    & 懊依曲(五首)                                           \\
                               & 谷雨后一日抵里（四选一）                           &                                                    &                                                        \\
    林纾                       & 叔节先生善余，……作此调叔节，并以自调               &                                                    &                                                        \\
    朱铭盘                     & 生日                                               &                                                    &                                                        \\
                               &                                                    &                                                    & 赠丘履平                                               \\
    释敬安                     &                                                    & 梦洞庭                                             &                                                        \\
                               &                                                    &                                                    & 梅痴子乞陈师曾为白梅写影，属赞三首(选一)               \\
    张骞                       &                                                    &                                                    & 屡出                                                   \\
                               &                                                    &                                                    & 病起                                                   \\
                               &                                                    &                                                    & 为黄仲弢编修绍箕题《龙女图》                           \\
    陈三立                     &                                                    & 短歌寄杨叔玫，时杨为江西巡抚令入红十字会观日俄战局 &                                                        \\
                               &                                                    & 晓抵九江作                                         &                                                        \\
                               &                                                    & 哭次申                                             &                                                        \\
                               &                                                    & 十一月十四夜发南昌月江舟行（四首选一）             &                                                        \\
                               &                                                    & 寄调伯弢高邮榷舍（二首选一）                       &                                                        \\
                               & 七月十四夜，结客十数辈泛东湖看月                   &                                                    &                                                        \\
                               & 遣怀                                               &                                                    &                                                        \\
                               & 伤邹沅颿                                           &                                                    &                                                        \\
                               &                                                    &                                                    & 江行杂感五首(选二)                                     \\
                               &                                                    &                                                    & 由沪还金陵散原别墅杂诗(五首选一)                       \\
                               &                                                    &                                                    & 野望                                                   \\
                               &                                                    &                                                    & 月夜                                                   \\
                               &                                                    &                                                    & 肯堂为我录其甲午客天津中秋玩月之作……用前韵奉报         \\
                               &                                                    &                                                    & 正月十九日园望                                         \\
                               &                                                    &                                                    & 春晴步后园晚望(五首选一)                               \\
                               &                                                    &                                                    & 渡湖至吴城(二首选一)                                   \\
    王树枏                     &                                                    &                                                    & 秋园居                                                 \\
    宋育仁                     &                                                    &                                                    & 甲午感事三首                                           \\
    范当世                     &                                                    & 吾所植荷既开尽而风雨频至坐见其萎谢慰别以诗         &                                                        \\
                               &                                                    & 大桥墓下                                           &                                                        \\
                               &                                                    & 过泰山下                                           &                                                        \\
                               & 轮船所谓买办唐姓者，余识之且二十年矣，见之而叹     &                                                    &                                                        \\
                               & 大桥墓下                                           &                                                    &                                                        \\
                               &                                                    &                                                    & 南康城下作                                             \\
                               &                                                    &                                                    & 守风至六七日之久，夜不复成寐，百虑交至，起眺书怀       \\
                               &                                                    &                                                    & 果然                                                   \\
                               &                                                    &                                                    & 暮春金陵城北见桃李花有感                               \\
    严复                       & 写怀                                               &                                                    &                                                        \\
                               &                                                    &                                                    & 赠畏庐                                                 \\
                               &                                                    &                                                    & 十三夜月                                               \\
                               &                                                    &                                                    & 古意                                                   \\
    程颂万                     &                                                    &                                                    & 迟明看荷呈恪士                                         \\
                               &                                                    &                                                    & 湖上草堂杂诗四首(选二)                                 \\
    王毓菁                     &                                                    &                                                    & 苔                                                     \\
    文廷式                     &                                                    & 夜坐向晓                                           &                                                        \\
                               & 辛丑元日试笔（二首）                               &                                                    &                                                        \\
                               & 自七月至十月，有感而作四首（选一）                 &                                                    &                                                        \\
                               &                                                    &                                                    & 谈仙诗(有序)                                           \\
                               &                                                    &                                                    & 拟古宫词(二十四首选一)                                 \\
    杨锐                       & 过卢沟                                             &                                                    &                                                        \\
    易顺鼎                     & 偕伯严云锦亭夜坐，有怀梁节庵去夏同游               &                                                    &                                                        \\
                               & 平望                                               &                                                    &                                                        \\
                               & 自襄樊至西安，道中绝句                             &                                                    &                                                        \\
                               &                                                    &                                                    & 天童山中月夜独坐二首                                   \\
    潘飞声                     &                                                    &                                                    & 罗浮纪游二首                                           \\
    沈瑜庆                     & 百花洲楼成，赠伯严                                 &                                                    &                                                        \\
    林鹤年                     & 丁酉四月初七日，厦口东望台、澎，泣而有赋           &                                                    &                                                        \\
    林寿图                     & 京口阻雨                                           &                                                    &                                                        \\
    陈衍                       &                                                    & 冬述四首示子培（四首选一）                         & ✓                                                      \\
    沈汝瑾                     &                                                    &                                                    & 食蟹                                                   \\
    刘光第                     & 屯海戍                                             &                                                    &                                                        \\
                               &                                                    & 华严顶                                             &                                                        \\
                               &                                                    &                                                    & 望峨眉山                                               \\
                               &                                                    &                                                    & 峨眉最高顶(在锡瓦殿后)                                 \\
                               &                                                    &                                                    & 夜灯                                                   \\
    康有为                     &                                                    & 秋登越王台                                         &                                                        \\
                               & 游花嫩冈，谒华盛顿墓宅                             &                                                    &                                                        \\
                               & 戊戌八月纪变，八首（选一）                         &                                                    & ✓                                                      \\
                               &                                                    &                                                    & 出都留别诸公(五首选三)                                 \\
    俞明震                     &                                                    & 宿凉州                                             &                                                        \\
                               &                                                    & 焦山松寥阁夜坐                                     &                                                        \\
                               &                                                    & 游西溪归泛舟湖上晚景奇绝和散原作                   &                                                        \\
                               &                                                    &                                                    & 同伯严后湖观荷                                         \\
                               &                                                    &                                                    & 湖庄晓起                                               \\
                               &                                                    &                                                    & 登剑门峰拂水桥                                         \\
                               &                                                    &                                                    & 天竺                                                   \\
                               &                                                    &                                                    & 春寒春楼写意                                           \\
                               &                                                    &                                                    & 夜雨待萧稚泉不至                                       \\
                               &                                                    &                                                    & 章江晚洎                                               \\
                               &                                                    &                                                    & 病起夜坐                                               \\
                               &                                                    &                                                    & 雨后湖楼晓起                                           \\
    夏曾佑                     &                                                    &                                                    & 别任公(二首选一)                                       \\
                               &                                                    &                                                    & 已亥秋别天津有感寄怀严蒋陈诸故人(四首选一)             \\
                               &                                                    &                                                    & 舟过大沽望炮台二首                                     \\
                               &                                                    &                                                    & 无题(二十六首选一)                                     \\
    萧蜕                       &                                                    &                                                    & 赠金东雷                                               \\
                               &                                                    &                                                    & 游小云栖                                               \\
    陈诗                       &                                                    &                                                    & 哭五弟子修诗并序                                       \\
                               &                                                    &                                                    & 乙巳冬日，与伯严先生饮酒垆，即送归秣陵                 \\
    何振岱                     &                                                    &                                                    & 鹤涧小坐                                               \\
                               &                                                    &                                                    & 理安寺                                                 \\
                               &                                                    &                                                    & 孤山独坐雪意甚足                                       \\
                               &                                                    &                                                    & 人杭舟中杂咏(二首)                                     \\
                               &                                                    &                                                    & 铁塔                                                   \\
                               &                                                    &                                                    & 高昌庙独步看柳                                         \\
                               &                                                    &                                                    & 舟夜                                                   \\
    李希圣                     & ✓                                                  & 西苑                                               &                                                        \\
                               &                                                    &                                                    & 帝子                                                   \\
    丘逢甲                     &                                                    & 题兰史《罗浮纪游图》                               &                                                        \\
                               &                                                    & 铁汉楼怀古                                         &                                                        \\
                               &                                                    & 西贡杂诗（十首选一）                               &                                                        \\
                               &                                                    & 韩江有感                                           &                                                        \\
                               &                                                    &                                                    & 海军衙门歌同温慕柳同年作                               \\
                               &                                                    &                                                    & 寄怀维卿师桂林(八首选一)                               \\
                               &                                                    &                                                    & 秋溪即目(一首)                                         \\
                               &                                                    &                                                    & 山村即目(选二)                                         \\
                               &                                                    &                                                    & 春日杂诗(四首选一)                                     \\
    姚永概                     &                                                    &                                                    & 偶题                                                   \\
    曾习经                     &                                                    &                                                    & 袁珏山所藏潘莲堂焦山图                                 \\
    黄人                       &                                                    &                                                    & 题长吉集                                               \\
                               &                                                    &                                                    & 萤                                                     \\
    罗惇曧                     &                                                    &                                                    & 题罗两峰《鬼趣图》                                     \\
                               &                                                    &                                                    & 登清凉山                                               \\
                               &                                                    &                                                    & 万松寺待晓                                             \\
                               &                                                    &                                                    & 上盘顶云罩寺暮不达而返                                 \\
                               &                                                    &                                                    & 香山雨香岩杂诗                                         \\
    谭嗣同                     &                                                    & 题宋徽宗画鹰                                       &                                                        \\
                               & 金陵听说法，三首（选一）                           &                                                    &                                                        \\
                               &                                                    &                                                    & 夜泊                                                   \\
    蒋智由                     &                                                    & 久思                                               &                                                        \\
    曾广钧                     &                                                    & 携眷登南岳观音岩作                                 &                                                        \\
    王允皙                     &                                                    & 梅花                                               &                                                        \\
                               &                                                    &                                                    & 于役书见                                               \\
    赵熙                       &                                                    & 龙门峡道中                                         &                                                        \\
                               & 上任父                                             &                                                    &                                                        \\
                               &                                                    &                                                    & 嘉定舟中                                               \\
                               &                                                    &                                                    & 秋夜                                                   \\
    丁惠康                     &                                                    &                                                    & 闻胶州近事二首                                         \\
    杨增荦                     & 讯伯严（二首选一）                                 &                                                    &                                                        \\
    章炳麟                     &                                                    & 狱中赠邹容                                         &                                                        \\
                               &                                                    &                                                    & 黑龙潭                                                 \\
    陈去病                     &                                                    &                                                    & 焦山中流遇急湍                                         \\
                               &                                                    &                                                    & 自厦门泛海登鼓浪屿有感                                 \\
                               &                                                    &                                                    & 癸卯除夕别上海，甲辰元旦……(八首选一)                   \\
    吴保初                     &                                                    & 乙巳游日本绝句                                     &                                                        \\
    梁启超                     &                                                    & 太平洋遇雨                                         &                                                        \\
                               & 奉怀南海先生星加坡，兼请东渡（二首选一）           &                                                    &                                                        \\
                               &                                                    &                                                    & 感秋杂诗(六首选三)                                     \\
                               &                                                    &                                                    & 庚戌岁暮感怀(六首选一)                                 \\
    黄节                       & 闭门                                               &                                                    &                                                        \\
                               &                                                    &                                                    & 宴集桃李花下，兴言边患，夜分不寐                       \\
                               &                                                    &                                                    & 报宾虹寄画                                             \\
                               &                                                    &                                                    & 岁暮示秋枚                                             \\
                               &                                                    &                                                    & 二月十二日过新汀屈翁山先生故里……谨题二首(选一)         \\
                               &                                                    &                                                    & 七月初六赴沪，……月中有怀曼殊                           \\
                               &                                                    &                                                    & 南归至沪寄京邸旧游                                     \\
                               &                                                    &                                                    & 送贞壮南归                                             \\
                               &                                                    &                                                    & 社园茗座迟瘿公不至                                     \\
    许承尧                     &                                                    & 言天                                               &                                                        \\
    金天羽                     &                                                    & 车中望居庸关放歌                                   &                                                        \\
                               &                                                    &                                                    & 罱泥船                                                 \\
                               &                                                    &                                                    & 悯农                                                   \\
                               &                                                    &                                                    & 看蚕娘                                                 \\
                               &                                                    &                                                    & 嵩山高                                                 \\
    林旭                       &                                                    & 马房沟                                             &                                                        \\
                               &                                                    & 无题                                               &                                                        \\
                               & 外舅《哀余皇》诗题后                               &                                                    &                                                        \\
                               &                                                    &                                                    & 张园梅花                                               \\
                               &                                                    &                                                    & 直夜                                                   \\
                               &                                                    &                                                    & 上海胡家闸茶楼                                         \\
    秋瑾                       &                                                    & 感愤                                               &                                                        \\
                               &                                                    & 咏燕                                               &                                                        \\
                               &                                                    & 秋海棠                                             &                                                        \\
                               &                                                    &                                                    & 秋风曲                                                 \\
                               &                                                    &                                                    & 感事                                                   \\
                               &                                                    &                                                    & 杞人忧                                                 \\
                               &                                                    &                                                    & 菊                                                     \\
                               &                                                    &                                                    & 对酒                                                   \\
    诸宗元                     &                                                    &                                                    & 夜过海藏楼，归纪所语简太夷，并示拔可                   \\
                               &                                                    &                                                    & 净慈寺井                                               \\
                               &                                                    &                                                    & 夜游约同去病                                           \\
                               &                                                    &                                                    & 徐州                                                   \\
                               &                                                    &                                                    & 雨中夜发上海，晓晴达金陵，复渡江趋浦口，舟中感纪       \\
                               &                                                    &                                                    & 哀寒碧                                                 \\
    周达                       &                                                    &                                                    & 题陈叔通所藏江弢叔手书诗卷                             \\
    庞树柏                     &                                                    &                                                    & 己酉十月朔，南社第一次雅集于虎溪张公祠，到者凡一十七人 \\
                               &                                                    &                                                    & 舟行西郭即景                                           \\
                               &                                                    &                                                    & 西泠杂咏八首(选四)                                     \\
    郁华                       &                                                    &                                                    & 马关                                                   \\
                               &                                                    &                                                    & 东京竹枝词(六首选一)                                   \\
    柳亚子                     &                                                    &                                                    & 哭宋遁初烈士                                           \\
                               &                                                    &                                                    & 吊鉴湖秋女士(选二)                                     \\
                               &                                                    &                                                    & 寄题西湖岳王冢同慧云作                                 \\
                               &                                                    &                                                    & 消寒                                                   \\
                               &                                                    &                                                    & 海上赠刘三                                             \\
    杨圻                       &                                                    & 题五洲地图                                         &                                                        \\
                               &                                                    & 京口遇范肯堂先生                                   &                                                        \\
                               &                                                    & 泰山玉皇顶                                         &                                                        \\
                               &                                                    &                                                    & 得幼儿丰祚贞祚家书                                     \\
    夏敬观                     &                                                    & 云栖寺竹径                                         &                                                        \\
                               &                                                    &                                                    & 湖上三首(选一)                                         \\
                               &                                                    &                                                    & 松径步月                                               \\
                               &                                                    &                                                    & 渡淮                                                   \\
                               &                                                    &                                                    & 雨过二首                                               \\
                               &                                                    &                                                    & 江夜                                                   \\
    陈衡恪                     &                                                    &                                                    & 月下写怀                                               \\
                               &                                                    &                                                    & 忆石湖旧游                                             \\
    林黻桢                     &                                                    &                                                    & 虎丘晚眺                                               \\
                               &                                                    &                                                    & 十三夜月                                               \\
    高旭                       &                                                    &                                                    & 读谭壮飞先生传感赋                                     \\
                               &                                                    &                                                    & 黄海舟中作                                             \\
    王国维                     &                                                    &                                                    & 读史二十首(选二)                                       \\
    于右任                     &                                                    &                                                    & 雨花台                                                 \\
                               &                                                    &                                                    & 题于鹤九画                                             \\
    宁调元                     &                                                    &                                                    & 从军行                                                 \\
    李宣龚                     &                                                    &                                                    & 徐州道中                                               \\
                               &                                                    &                                                    & 夜坐示贞壮并寄映庵江南                                 \\
                               &                                                    &                                                    & 焦山枕江阁同沤尹丈夜坐                                 \\
    汪荣宝                     &                                                    &                                                    & 故国                                                   \\
    马君武                     &                                                    &                                                    & 去国辞(五首选一)                                       \\
    傅熊湘                     &                                                    &                                                    & 大汉报社次和默君女士赠别韵(二首选一)                   \\
    周实                       &                                                    &                                                    & 岁除前一日作                                           \\
    柳景行                     &                                                    &                                                    & 畏暑人湖，得行严京国书却寄                             \\
    胡朝梁                     &                                                    &                                                    & 岁暮杂诗(五首选一)                                     \\
    陈曾寿                     &                                                    & 游仙                                               &                                                        \\
                               &                                                    & 观瀑亭                                             &                                                        \\
                               &                                                    & 湖斋坐雨                                           &                                                        \\
                               &                                                    & 湖上杂诗（七首选二）                               &                                                        \\
                               &                                                    &                                                    & 南湖晦夜寄怀散原先生(四首选一)                         \\
                               &                                                    &                                                    & 壬子二月同恪士梅庵至西湖刘氏花园                       \\
    宋教仁                     &                                                    &                                                    & 秋晓                                                   \\
    苏曼殊                     &                                                    & 本事诗十章（选一）                                 &                                                        \\
                               &                                                    &                                                    & 耶婆提病中，末公见示新作，伏枕奉答，兼呈旷处士         \\
                               &                                                    &                                                    & 简法忍                                                 \\
                               &                                                    &                                                    & 久欲南归罗浮不果……聊书所怀……                           \\
    林学衡                     &                                                    &                                                    & 南河沿                                                 \\
    林文                       &                                                    & 春望                                               &                                                        \\
    孙景贤                     &                                                    & 抵浦口（二首）                                     &                                                        \\
    龙文彬                     & 偶书                                               &                                                    &                                                        \\
    唐烜                       & 戊戌纪事八十韵                                     &                                                    &                                                        \\
    陈沣                       & 感旧                                               &                                                    &                                                        \\
    姚椿                       & 月夜望君山                                         &                                                    &                                                        \\
                               & 客中除夕                                           &                                                    &                                                        \\
    天培                       & 中秋，嶰筠尚书招余及关滋圃军门                     &                                                    &                                                        \\
    林则徐                     & 饮沙角炮台，眺月有作                               &                                                    &                                                        \\
                               & 次韵答宗涤楼（稷辰）赠行（二首）                   &                                                    &                                                        \\
    程恩泽                     & 西平道中                                           &                                                    &                                                        \\
                               & 沉阴                                               &                                                    &                                                        \\
    吴振棫                     & 三叹                                               &                                                    &                                                        \\
                               & 云里                                               &                                                    &                                                        \\
    梅曾亮                     & 戏书                                               &                                                    &                                                        \\
                               & 监利王子寿去刑部主政归，作诗寄之                   &                                                    &                                                        \\
                               & 偶成                                               &                                                    &                                                        \\
    徐荣                       & 西岭道中                                           &                                                    &                                                        \\
    宗稷辰                     & 舟行晚望                                           &                                                    &                                                        \\
                               & 途中书见（二首）                                   &                                                    &                                                        \\
    祁寯藻                     & 灰槌，相传始皇焚书处（二首选一）                   &                                                    &                                                        \\
                               & 送三侄世龄还文水学舍                               &                                                    &                                                        \\
                               & 题壹斋师饯书图（二首选一）                         &                                                    &                                                        \\
                               & 立夏后一日，长椿寺牡丹初开,漫题                    &                                                    &                                                        \\
    张际亮                     & 酒中书扇，赠林二观成需然                           &                                                    &                                                        \\
    黄爵滋                     & 抚州行，寄太守郭松厓前辈                           &                                                    &                                                        \\
                               & 不见                                               &                                                    &                                                        \\
    彭蕴章                     & 远行篇                                             &                                                    &                                                        \\
    李联琇                     & 悔昔篇                                             &                                                    &                                                        \\
    蒋敦复                     & 独行海上，作歌告哀，言之无罪，惜当事者终竟不闻也   &                                                    &                                                        \\
    徐子苓                     & 中秋后十日夜，书啸和尚水灾诗后                     &                                                    &                                                        \\
                               &                                                    &                                                    & 腊月廿四日，遣郑仆往周云先家迎吴四引之二首(选一)       \\
    王尚辰                     &                                                    &                                                    & 荻港吊黄靖南                                           \\
                               &                                                    &                                                    & 花下                                                   \\
  \end{longtblr}
\end{center}
% \end{landscape}


%%% Local Variables:
%%% TeX-engine: xetex
%%% TeX-master: "Qing_Poem_300_Compare"
%%% End:

% LocalWords:  utf

%% -*- coding: utf-8 -*-

\section{诗人简介}

\newcounter{myterm}
\begin{description}[
  before={\setcounter{myterm}{0}}, % 关键：每次进入这个列表时，把序号清零
  font={\stepcounter{myterm}\arabic{myterm}.~\bfseries}, % 关键：每次遇到 \item，数字加1并输出，然后保持名词粗体
  labelsep=0.5em,
  leftmargin=3em % 稍微增加一点左侧缩进，给序号留出空间
]
  \item [钱仲联] （1908年—2003年），初名萼孙，字仲联，号梦苕，以字行，江苏常熟人，原籍浙江省吴兴县，中国文学研究家，任教于大夏大学、无锡国学专修学校、南京中央大学、苏州大学等校。有“昆仑万象一代诗豪”、“清诗功臣”之美誉。
  \item [刘世南] （1923年10月-2021年8月1日），江西吉安人，文史学者、诗人。曾任江西师范大学文学院副教授、硕士生导师，1989年退休，2006年补评为教授，任《豫章丛书》点校首席学术顾问及江西省古籍整理中心组成员。\par
        幼年随父研习古籍，高一辍学后长期在吉安、新建等地任中学语文教师。1979年因发表《谈古文的标点、注释和翻译》等论文引发学界关注，同年调入江西师范大学任教。代表专著《清诗流派史》由台北文津出版社和人民文学出版社先后出版，被学界视为清诗研究经典著作，另著有《清诗三百首详注》《黄遵宪诗选注》《清文选》《谷梁传直解》《古文观止译注》《在学术殿堂外》《大螺居诗存》《大螺居诗文存》《师友偶记》等。旧体诗作入选《海岳弦歌集》等选集，并为高校撰拟楹联铭文。2021年8月1日在南昌逝世，享年99岁。
  \item [\hypertarget{XingFang}{邢昉}] （1590年—1653年），字孟贞，一字石湖。江苏高淳人。\par

        曾祖邢符，祖邢尚宾，父邢一淳，科场俱不得意。明末诸生，屡试不中，第六次参加乡试，主考官称其文笔太狂﹐愤而作《太狂篇》，从此绝意仕途。与方文友好，两人有诗往来。崇祯十年丁丑（1637年），应杨文骢之请赴华亭幕府。崇祯十三年庚辰（1640年），邢昉、方文等人于金陵结社。崇祯十四年辛巳（1641年），方文赴庐州投靠蔡如蘅，邢昉有诗相送。崇祯十六年癸未（1643年），邢昉旅居宣城。崇祯十七年甲申（1644年）八月，母赵孺人逝。\par
        邢昉是明代著名诗人，终身布衣，为诗清真古澹，以《故宫燕》最著名，王士禛在《渔洋诗话》独推邢昉“韦、柳门庭人”及“布衣诗人第一”。吴敬梓曾拜访其石臼湖畔故居，谒邢昉墓。著有《宛游草》、《石臼集》。
\end{description}


%%% Local Variables:
%%% TeX-engine: xetex
%%% TeX-master: "Qing_Poem_300_Compare"
%%% End:


% \input{contents/part000}

\end{document}

% Local Variables:
% TeX-engine: xetex
% TeX-master: t
% End:

% LocalWords:  xetex AutoFallBack AutoFakeBold PunctStyle kaiming CJKmath
% LocalWords:  CheckFullRight AllowBreakBetweenPuncts ctexset
